\section{Conclusion}
This report investigated the availability threat posed by Sponge Attacks to AI-based Smart Grid systems. We identified that static RNNs (DeepAR) and Foundation Models (Chronos-v2) are robust to \textit{latency}-based sponge attacks due to their fixed computational graphs, while dynamic architectures like ACT-LSTM are critically vulnerable. However, all architectures—including static ones—are susceptible to the bit-flip energy attack. Our Proof-of-Concept demonstrated that a black-box Genetic Algorithm could induce a 53\% latency spike and 70\% increase in ponder depth in dynamic models by exploiting the logical decision boundaries of the halting unit.

Furthermore, our bit-flip oracle analysis revealed that dynamic power dissipation—and therefore energy consumption—scales proportionally with model size across \textit{all} architectures, not just dynamic ones. Foundation models like Chronos-v2, with their large embedding layers, produce $470\times$ more bit-flip transitions than compact models like ACT-LSTM, making them inherently more vulnerable to switching-activity energy attacks. This finding suggests that the trend towards larger models in Smart Grid AI inadvertently amplifies a second, architecture-agnostic energy attack surface.

As Smart Grid operators increasingly look towards dynamic and adaptive AI for efficiency, the attack surface for Denial-of-Service expands. We conclude that \textit{availability defenses}—such as strict compute timeouts, architectural distillation, and hardware-aware input validation—are mandatory requirements for deploying deep learning in time-critical energy infrastructure.
